\section{Research Directions y Q\&A: lo que todavía no se sabe}

La clase no cerró con «scale solves everything», sino con una lista: uncertainty, inference-time learning, adaptive thinking, planning y quantization. La sesión de Q\&A añadió además condiciones limitantes como non-autoregressive ordering, language world knowledge, generalization, multi-agent coordination, modularity y gradient descent. Estos spoken nodes son la parte de esta clase donde más se concentra la teacher voice.

\subsection{Cinco Research Directions públicas}

Primero, el modelo debería saber lo que no sabe; segundo, si el modelo puede aprender un skill nuevo en inference time, sin actualizar sus parámetros cada vez; tercero, si puede asignar adaptive compute según la dificultad del problema; cuarto, si el reasoning/planning puede resolverse con un modelo pequeño, un external planner o una mejor representation; quinto, si una precision más baja puede seguir reduciendo el memory traffic y aumentando el matrix throughput.

\lecturefigure{slide-43-research-directions.jpg}{Incertidumbre, inference-time learning, adaptive thinking, planning y quantization}{Diapositiva V043 recuperada de la grabación oficial; 01:00:51}

\begin{knowledgebox}{Lectura de la figura: las cinco direcciones abarcan model, algorithm y system}
Knowing what you do not know pertenece a calibration/evaluation; learning skills at inference time pertenece a memory/adaptation; adaptive thinking y planning pertenecen a compute allocation; quantization pertenece al numerical/system co-design. No pueden validarse con un solo indicador de «modelo más grande».
\end{knowledgebox}

\teachervoice{El ponente observa que un diffusion model puede mejorar su calidad con más compute de muestreo, mientras que los modelos de lenguaje carecían entonces de un adaptive-compute knob igual de claro. Plantea la pregunta: ¿se puede hacer que un modelo pequeño piense más en los problemas difíciles y se conecte a un external planner, para acercarse al reasoning de un modelo grande, en lugar de pagar el costo máximo por cada token?}

\subsection{Non-autoregressive Decoding: lo que realmente falta es un Order Oracle}

La sesión de Q\&A explicó con más detalle la parallel generation. Si el latent/context ya generado basta para que el siguiente grupo de tokens sea condicionalmente independiente, ese grupo puede producirse en paralelo; la dificultad está en aprender ese grouping/order. Una factorization por grupos es:
\[
p(y\mid x)=\prod_{g=1}^{G}\prod_{i\in S_g}p\!\left(y_i\mid x,y_{S_{<g}}\right),
\]
donde $S_g$ es el $g$-ésimo token group que puede generarse en paralelo y $S_{<g}$ son los groups generados antes. Si un oracle proporciona los groups correctos, el problema se simplifica mucho; en la práctica, el modelo además tiene que descubrir la conditional independence y manejar distribuciones de salida multimodales.

\teachervoice{El ponente dice que, si alguien proporcionara un oracle ordering del tipo «primero genera estas tres palabras y luego esas dos», la non-autoregressive generation tendría más posibilidades. Intentaron modelar las dependencias con una vector-quantized latent sequence; funcionó en translation, pero requería distillation para reducir la entropy de los datos y, al final, su practical impact fue menor que el de speculative decoding.}

\subsection{Language-only World Model: insuficiente, pero más útil de lo que se intuye}

La clase no aceptó que «el texto permite aprender todo el mundo físico», pero tampoco negó su usefulness. Los modelos de lenguaje grandes absorben del texto una gran cantidad de conocimiento sobre objetos, procedimientos y relaciones sociales, y pueden servir como prior de alto nivel para un robotics planner; la percepción, el control y la retroalimentación del entorno siguen a cargo de otros sistemas. Lo esencial es distinguir knowledge coverage de action grounding.

\begin{warningbox}{Poder planificar no equivale a poder controlar en lazo cerrado}
Que un modelo de lenguaje enuncie pasos razonables solo demuestra que puede organizar sus priors textuales; el robot además necesita reconocer los objetos presentes, estimar el estado, cumplir con la dinámica, ejecutar de forma segura y recuperarse de los fallos. El world knowledge es planner input, no physical verification.
\end{warningbox}

\teachervoice{La postura de Vaswani es pragmática: quizá language alone sea incompleto, pero los robots ya pueden usar los modelos grandes como planner, lo que muestra que el texto contiene más información sobre el mundo de lo que muchos esperaban. Ni siquiera hemos aprovechado por completo la usefulness de los modelos actuales, así que una insuficiencia de fondo no basta para negar su valor de ingeniería actual.}

\subsection{Generalization: las capacidades nuevas pueden surgir de una Recombination a gran escala}

Si los datos de entrenamiento tienen una cobertura muy amplia, el límite OOD tradicional se vuelve difícil de definir. Puede que el modelo nunca haya visto un ejemplo exacto de «describir esta Stanford lecture con la métrica de Chaucer», pero sí ha visto clases, a Chaucer, métrica y reescrituras, y puede combinarlos en una salida nueva. Esto es compositional generalization útil, aunque sus piezas de información atómicas provengan de la distribución de entrenamiento.

\begin{knowledgebox}{No hay que reducir la Recombination a «solo memoria»}
Una evaluación debe distinguir entre exact memorization, template interpolation, systematic composition y algorithmic extrapolation. La recombination a gran escala puede generar valor real, pero no garantiza una extrapolación arbitraria en longitud, números, reglas o espacio de estados; cada benchmark mide una generalization distinta.
\end{knowledgebox}

\subsection{Multi-agent Systems: Decomposition, Coordination, Verification}

Que varios agents superen a un solo modelo grande no depende del número de roles. El ponente plantea tres preguntas básicas: cómo descomponer el objetivo según el skill de cada agent, cómo coordinar el proceso con estado compartido y dependencias, y cómo verificar el resultado. Sin verification, más agents solo producen más texto y una failure propagation más compleja.

\begin{importantbox}{Los tres requisitos mínimos de validación de un sistema Multi-agent}
\textbf{Goal decomposition}: dividir la tarea en subgoals verificables; \textbf{coordination}: manejar el orden, los recursos, los conflictos y la memory compartida; \textbf{verification}: comprobar los resultados con pruebas, estado del entorno o evidencia externa. Si falta cualquiera de los tres, el llamado swarm no es más que muestreo en paralelo.
\end{importantbox}

\teachervoice{El ponente dice medio en broma que «los mejores agents son las neurons», porque se actualizan de forma coordinada a través del gradient. Para los sistemas multi-agent explícitos no ofrece una arquitectura universal, sino que reconoce abiertamente que decomposition, coordination y verification siguen en una etapa temprana.}

\subsection{Modularity, MoE y la «Architecture al servicio del Gradient Descent»}

Los module diseñados por personas suelen fracasar por las dificultades de routing, credit assignment y optimización conjunta. Mixture of Experts (MoE) ofrece learned routing, de modo que cierta specialization surge dentro del end-to-end gradient descent; pero la expert identity no equivale a un módulo cognitivo completo, y la responsabilidad puede seguir distribuida entre la attention, el residual stream y varios experts.

\teachervoice{Vaswani dice que a veces siente que «solo estamos construyendo architecture para servir al gradient descent». La frase no desprecia la arquitectura; señala que la entrenabilidad es una restricción estricta: por elegante que sea una división en módulos, si el credit no puede propagarse y el routing no puede aprenderse, difícilmente sustituirá a una red más homogénea pero optimizable.}

\subsection{Resumen de la sección}

Las research directions abarcan calibration, adaptation, compute allocation, planning y numerical efficiency. La sesión de Q\&A muestra además que la parallel generation requiere una estructura de order/independence, que el language world knowledge es útil pero no equivale a grounding, que los sistemas multi-agent requieren verification y que la modularity debe subordinarse a la entrenabilidad.
